\documentclass[11pt]{article}

%%% Useful packages
\usepackage{amsfonts}
\usepackage{amssymb}
\usepackage{amsmath}
\usepackage{amsthm}

%%% Sets page margins, overriding 11pt article defaults
\setlength{\topmargin}{0 in}
\setlength{\headheight}{0 in}
\setlength{\headsep}{0 in}
\setlength{\oddsidemargin}{0in}
\setlength{\textwidth}{6.5 in}
\setlength{\textheight}{9 in}

%%% Useful for layout
\newcommand{\VF}{\vspace*{\fill}}
\newcommand{\HF}{\hspace*{\fill}}

%%% New commands used


\begin{document}

{\noindent\Large\textbf{Homework 3 due September 18}}

\vspace{.25in}

{\large
\noindent
\textbf{Name:} \smallskip \\
\textbf{GTID:} \smallskip \\
%%% List anyone with whom you discussed the problems
\textbf{Collaborators:} \smallskip \\
%%% List all resources used OTHER than the textbook, lecture notes, 
%%% webwork problems, ICA questions, and previous homework assignments.
\textbf{Outside resources:} \smallskip \\
}

\pagebreak

\smallskip

\noindent
In addition to results already proved in class or on previous ICA, Wbwk, 
and Hmwk assignments,
you may use any result proved in Sections 4.4, 4.5, and 4.6 of the textbook.

\smallskip

\noindent
You may \emph{not} simply state ``Without loss of generality'' or any 
other similar shortcut.
Instead, after proving the first case completely, 
briefly explain why the other case(s) do not need to be 
considered in detail.

\bigskip

\noindent
Let $A$, $B$, $C$, and $D$ be subsets of a universal set $\mathcal{U}$.

\smallskip

\begin{enumerate}

\item Prove:
\begin{enumerate}
\item $(A \cup B) = (A \cap B) \cup (A \setminus B) \cup (B \setminus A)$.  
\item $(A \cap B)$, $(A \setminus B)$, and $(B \setminus A)$ are all 
pairwise disjoint.
\end{enumerate}
\emph{This proves that the three subsets are a \emph{partition} of 
$(A \cup B)$.  
Partitions will be very important to understanding equivalence relations.}

\item Prove or disprove:
\begin{enumerate}
\item $(A \times B) \cap (C \times D) \subseteq (A \cap C) \times (B \cap D)$.
\item $(A \cap C) \times (B \cap D) \subseteq (A \times B) \cap (C \times D)$.
\item $(A \times B) \cup (C \times D) \subseteq (A \cup C) \times (B \cup D)$.
\item $(A \cup C) \times (B \cup D) \subseteq (A \times B) \cup (C \times D)$.
\end{enumerate}
\emph{Likewise, Cartestian product of sets (above) and the power set of a 
set (below) will be very important in the next part of the course.}

\item Prove or disprove:
\begin{enumerate}
\item
\( {\mathcal P}(A\cap B)\subseteq {\mathcal P}(A) \cap {\mathcal P}(B) \)
\item
\( {\mathcal P}(A) \cap {\mathcal P}(B)\subseteq {\mathcal P}(A\cap B) \)
\item
\( {\mathcal P}(A\cup B)\subseteq {\mathcal P}(A) \cup {\mathcal P}(B) \)
\item
\( {\mathcal P}(A) \cup {\mathcal P}(B)\subseteq {\mathcal P}(A\cup B) \)
\end{enumerate}

\end{enumerate}

\end{document}

